% !TeX root = RJwrapper.tex
\title{couplr: Optimal Matching with Verifiable Certificates and Sparse Edge Generation}


\author{by Gilles Colling}

\maketitle

\abstract{%
Matching two groups of units so that the total cost of the chosen pairs is as small as possible is the linear assignment problem, which sits underneath causal-inference matching and experimental design. Two things bound what software built on it can offer. An answer arrives without a proof: a solver reports the state it terminated in, not a checkable statement about the matching. And the graph of admissible pairs is built before it is solved, so its memory sets the size ceiling. We present the couplr package, which addresses both. Every flow-representable matching design compiles into one internal flow model, so node potentials come back from all of them and can be checked against the linear-programming optimality conditions. On those potentials the package optimises over a restricted set of pairs: each unit starts with its nearest admissible partners, the omitted pairs are priced against the duals the sparse solve returns, and pairs enter until none prices in, at which point the duals certify the result optimal for the complete problem. Matching 16,667 treated units to 33,333 controls holds 0.29 percent of the pairs and returns a certified optimum. Warm starting the same loop traces a matching frontier over a swept constraint at roughly half the cost of solving each point cold. A data-frame interface, nineteen assignment algorithms, balance diagnostics and cardinality matching share the same core. couplr is available on CRAN.
}

\section{Introduction}\label{introduction}

Many data-analysis tasks come down to the same question: given two sets of
objects and a cost attached to every possible pair, which pairing has the
smallest total cost? A researcher building an observational study pairs treated
units with controls that resemble them on measured covariates, a scheduler
assigns staff to shifts, an ecologist pairs surveyed plots across two time
points, and an image-processing pipeline maps the pixels of one picture onto the
pixels of another. In each case the mathematical object is the linear assignment
problem, and the quality of the answer depends on solving it to optimality
rather than greedily.

R has good software at both ends of this problem, in two separate groups of
packages. On the statistical side, \CRANpkg{MatchIt} \citep{MatchIt} provides a
common interface to several matching designs and delegates optimal pair and
full matching to \CRANpkg{optmatch} \citep{optmatch}, which formulates those
designs as minimum-cost flow problems \citep{HansenKlopfer2006}, reaches RELAX-IV
or one of four LEMON algorithms through its \texttt{solver} argument, and stores
sparse distances as an \texttt{InfinitySparseMatrix} whose non-finite entries are the
pairings it will not consider. \CRANpkg{Matching} \citep{Matching} provides
multivariate and propensity-score matching, with \texttt{GenMatch()} using genetic
search to optimise covariate weights. \CRANpkg{designmatch} \citep{designmatch}
formulates balance-constrained designs as integer programs.
\CRANpkg{rcbalance} \citep{rcbalance} matches under refined covariate balance
constraints \citep{Pimentel2015}, taking a distance structure that lists the
permissible controls of each treated unit and solving the resulting network as
a minimum-cost flow problem. \CRANpkg{quickmatch} \citep{quickmatch} aims at large
data sets with generalized full matching, which its own description calls
near-optimal. \CRANpkg{cobalt} \citep{cobalt} supplies balance diagnostics across
packages. What these expose is the matching design; the assignment problem
underneath is reached through a solver argument at most, and alternative optima
and a bottleneck objective sit outside all of them. Statements about other
packages here and in Table \ref{tab:capability} were assessed on 2026-08-09
against \CRANpkg{MatchIt} 4.7.2, \CRANpkg{optmatch} 0.10.8, \CRANpkg{Matching} 4.10-15 and
\CRANpkg{designmatch} 0.5.5, and against the reference manuals of
\CRANpkg{rcbalance} 1.8.8 and \CRANpkg{quickmatch} 0.2.3 on 2026-08-26.

On the algorithmic side, \CRANpkg{clue} \citep{clue} exposes the Hungarian method
through \texttt{solve\_LSAP()}, and \CRANpkg{lpSolve} \citep{lpSolve} solves the assignment
problem as a general linear program. These return an optimum for a cost matrix
the user has already built; covariate handling, calipers and balance tables sit
outside their scope.

A workflow that needs both halves therefore threads several packages by hand.
Matching on observed covariates instead of a propensity score, imposing
calipers on more than one variable at a time, blocking, and choosing the solver
deliberately are all available in R, each in a different package.

Two further limits are common to both groups, and they are the ones this
article is about. The first concerns what a result says about itself. A solve
returns a matching and a status, and the status records the state the solver
terminated in. Whether that matching is the optimal one is a separate question with a
checkable answer, since linear programming supplies optimality conditions a
third party can test on the returned solution. The general linear programming
interfaces return material for that test in their own terms: \CRANpkg{lpSolve}
returns a \texttt{duals} component from \texttt{lp()} and \texttt{lp.assign()} when \texttt{compute.sens}
is set, indexed by the rows and columns of the program as assembled, so reading
those numbers as assignment potentials and checking the conditions against a
cost matrix is left to the caller. Among the matching-specific interfaces
surveyed, none exposes assignment or flow potentials together with a function
that verifies optimality against them. \CRANpkg{optmatch} states bounds in
their place, an \texttt{exceedances} attribute and a solver tolerance, both described
with the capability table below.

The second concerns size. Sparse matching is a developed literature, and what
its methods have in common is that the sparse network is settled before the
matching is solved over it. \CRANpkg{rcbalance} takes a per-unit list of
permissible controls built from exact groups and a caliper. \citet{YuSilberRosenbaum2020} choose a propensity-score caliper by an iterative form
of Glover's algorithm and match optimally in the much sparser graph that
leaves, and \citet{YuRosenbaum2022} grade the potential pairings and admit lower
grades only as far as they are needed. In each case the rule deciding which pairs
exist is fixed in advance of the solve that uses them, and a network too sparse
drops pairings it should have considered. What is missing is a rule read off the
solution in hand, and a statement about the pairs it never built.

We present the \CRANpkg{couplr} package \citep{couplr}, which treats the assignment
problem itself as the organising object and builds the workflow around it. A
user who starts from a data frame with covariates and a treatment indicator,
and a user who starts from a cost matrix, enter through different functions and
reach the same solver core. We contribute five things to R:

\begin{enumerate}
\def\labelenumi{\arabic{enumi}.}
\tightlist
\item
  A common compilation layer. One-to-one, k-to-one, with replacement, blocked
  and full matching are the same internal flow model under different arc
  capacities, so what is written once serves all of them.
\item
  Certificates rather than reported status. \texttt{verify\_assignment()} and
  \texttt{verify\_flow()} check primal feasibility, dual feasibility, complementary
  slackness and objective equality on a solution, reporting each separately
  and independently of the solver that produced it, in exact arithmetic where
  the potentials allow and naming the arithmetic the verdict was reached in.
\item
  Dual-certified adaptive edge generation. Under
  \texttt{memory\_mode\ =\ "implicit"} the complete graph is never built: a sparse
  restricted problem is solved, the omitted pairs are priced against its
  duals, violating pairs are added, and the loop repeats until none prices in,
  at which point the duals certify optimality for the complete problem.
\item
  Warm-started matching frontiers. \texttt{match\_path()} sweeps one constraint over
  a sequence of values, resuming each point from the one before it, with a
  certificate per point.
\item
  Nineteen assignment algorithms written from scratch in C++ and reachable by
  name, covering augmenting-path, auction, cost-scaling and network-flow
  families, with k-best and bottleneck assignment, entropy-regularised
  transport, and cardinality matching that reports an optimality gap against
  the largest sample its constraints admit.
\end{enumerate}

The article is arranged for two kinds of reader. The assignment engine and the
matching interface built on it come first and stand on their own; the
optimisation architecture, its implementation and the measurements follow for a
reader interested in where the certificates and the sparse edge generation come
from.

Results were produced with couplr version 1.7.0.

\section{The core assignment engine}\label{the-core-assignment-engine}

\subsection{The assignment problem}\label{the-assignment-problem}

Let \(C \in \mathbb{R}^{n \times m}\) be a cost matrix whose entry \(C_{ij}\) is the
cost of pairing row \(i \in \{1, \dots, n\}\) with column \(j \in \{1, \dots, m\}\),
and take \(n \le m\); couplr transposes the problem internally when it is not and
maps the assignment back afterwards, so the user never sees the transpose.
Write \(X \in \{0,1\}^{n \times m}\) for the assignment matrix, with \(X_{ij} = 1\)
when row \(i\) is matched to column \(j\). The linear assignment problem is

\begin{equation}
\min_{X} \; \sum_{i=1}^{n} \sum_{j=1}^{m} C_{ij} X_{ij}
\quad \text{subject to} \quad
\sum_{j} X_{ij} = 1 \; \forall i, \qquad
\sum_{i} X_{ij} \le 1 \; \forall j .
\label{eq:lap}
\end{equation}

Every row is matched and every column is used at most once. The row constraint
is an equality: relaxing it to \(\sum_j X_{ij} \le 1\) would leave \(X = 0\)
optimal for any non-negative cost matrix, which is not the problem a matching
application poses.

Forbidden pairs are encoded as \(C_{ij} = \infty\), and
enough of them make \eqref{eq:lap} infeasible, since no assignment can then
reach every row. couplr solves a lexicographic relaxation in that case: the
largest number of matched rows first, then the smallest total cost among the
assignments attaining that number. The section on constraints below describes
how that problem is reached and what it costs.

The constraint matrix of \eqref{eq:lap} is totally unimodular, so the linear
programming relaxation has integral optima \citep{Birkhoff1946, Burkard1999} and the
combinatorial problem can be solved by primal-dual methods rather than by branch
and bound. The dual assigns a potential \(u_i\) to each row and \(v_j\) to each
column, and is

\begin{equation}
\max_{u, v} \; \sum_{i=1}^{n} u_i + \sum_{j=1}^{m} v_j
\quad \text{subject to} \quad
u_i + v_j \le C_{ij} \; \forall (i,j), \qquad
v_j \le 0 \; \forall j .
\label{eq:duals}
\end{equation}

The sign condition is what the column inequality in \eqref{eq:lap} costs: the
dual objective sums \(v\) over every column while an assignment pays only for the
ones it uses. When \(n = m\) none can be left out, so the column inequalities may
be replaced by equalities without changing the feasible set, and the equality
formulation leaves \(v\) unrestricted in sign; \texttt{verify\_assignment()} imposes
\(v_j \le 0\) only
when \(n < m\), since Jonker-Volgenant returns free-sign potentials on a square
problem and they certify optimality there.

Complementary slackness ties a pair of potentials to a particular \(X\):

\begin{equation}
X_{ij} = 1 \implies u_i + v_j = C_{ij} , \qquad
\textstyle\sum_{i} X_{ij} = 0 \implies v_j = 0 .
\label{eq:cs}
\end{equation}

The reduced cost \(C_{ij} - u_i - v_j\) is a lower bound on what forcing pair
\((i,j)\) would add to the total rather than the whole of it, because forcing one
pair can displace others.

\subsection{Verifiable optimality certificates}\label{verifiable-optimality-certificates}

Dual feasibility on its own says nothing about a particular \(X\). What
certifies optimality is \(X\) primal feasible, \((u, v)\) feasible for
\eqref{eq:duals}, and both halves of \eqref{eq:cs} holding; in exact arithmetic
those three settle the objectives as well, so objective equality is a reported
check rather than a condition of its own. \texttt{verify\_assignment()} reports the
four separately, so a failure names the check it failed on.

\begin{verbatim}
set.seed(2)
cost <- matrix(runif(120), nrow = 6, ncol = 20)
verify_assignment(assignment(cost), cost)
\end{verbatim}

\begin{verbatim}
#> Assignment certificate
#> ======================
#> 
#>   primal_feasible          TRUE 
#>     valid matching         TRUE 
#>     all rows matched       TRUE 
#>   dual_feasible            TRUE 
#>   complementary_slackness  TRUE 
#>     matched arcs tight     TRUE    (max slack 0.000e+00)
#>     unmatched columns free TRUE    (max |v_j| 0.000e+00)
#>   duality_gap              0.000000e+00
#>   certified_optimal        TRUE 
#>   arithmetic               exact, no tolerance
#> 
#>   primal objective  0.2821665586
#>   dual objective    0.2821665586
#>   matched           6 of 6 rows, 20 columns
\end{verbatim}

Two properties make this a check rather than a report. It is written against
the solution rather than the solver, taking an assignment, a cost matrix and a
pair of potentials and never consulting the code that produced them. And
optimal duals are shared by all optimal solutions of a linear program, so a
matching from one solver can be certified against potentials from another,
which is what puts the solvers returning no duals of their own inside its
scope.

Complementary slackness is where a plausible verifier goes wrong, and the
rectangular case is where it shows. Asserting the first half of \eqref{eq:cs}
alone, tightness on matched pairs, accepts solutions that are not optimal: in
our own prototyping a perfect, dual-feasible, matched-tight solution passed
such a check with its primal cost 0.4\% above the optimum, because one
unmatched column retained \(v_j < 0\) and contributed a spurious term to the
dual objective. \texttt{verify\_assignment()} asserts both halves and reports them as
\texttt{cs\_matched\_tight} and \texttt{cs\_unmatched\_free}.

The arithmetic the conditions are decided in is the other place where a
certificate can overstate what it knows. Each of them is the sign of
\(C_{ij} - u_i - v_j\), and every quantity in that expression is a double, which
is a rational number, so each sign has an exact answer. Deciding them within a band around zero instead makes the verdict a statement
about a neighbourhood of the instance rather than about the instance, and what
the band costs can be written down. A certificate at tolerance \(\epsilon\)
accepts potentials that violate dual feasibility by \(\epsilon\) on any pair, miss
tightness by \(\epsilon\) on any matched arc, and leave \(|v_j| \le \epsilon\) on any
unused column. Shifting the potentials down by \(\epsilon\) restores exact dual
feasibility and the sign condition, and bounds the matching's excess over the
optimum by \((n + 2m)\epsilon\); the supplementary material carries the
derivation. At the largest shape in the benchmark below and the default
\(\epsilon = 10^{-9}\), that is \(8 \times 10^{-5}\) in the cost units.

Setting the band to zero is not the fix, because the sign of the double
evaluation can itself be wrong: a reduced cost far smaller than the potentials
it is formed from can round to zero, and in the negative case that is a dual
infeasibility a zero-band check would accept. \texttt{verify\_assignment()} evaluates
the sign exactly instead: it splits each rounded addition into the sum and the
part rounding lost, giving an expansion of doubles equal to the difference in
the reals, and takes the sign of that expansion \citep{Shewchuk1997}. A
rounding-error bound decides which pairs the double evaluation already settles,
so the expansion runs near tightness rather than on all \(nm\).

Removing the band from the two conditions removes it from the conclusion.
Objective equality is then not a fourth measurement but a consequence: with
matched arcs tight and unmatched columns free, the primal cost is the sum of
\(u\) over matched rows plus the sum of \(v\) over used columns, which is the dual
objective exactly when every row is matched. The software reports the gap it
computes as an independent numerical cross-check on that reasoning.

An exact certificate states that the matching attains the minimum total cost
of the matrix as supplied, and is available when the potentials are exactly
optimal for it, which is a property of the arithmetic that produced them. Over
ten instances per cell, integer costs and uniform costs on the unit interval
gave one in every instance at 50, 200 and 800 rows square, at 200 rows against
600 columns, and for each of the eleven solvers we ran at 200 by 200.
Euclidean distances between normal clouds gave one in none: those potentials
miss exact tightness by a relative \(2 \times 10^{-16}\) to \(5 \times 10^{-15}\),
which is what rounding in the arithmetic that produced them costs, and there
the certificate decides the same conditions within a stated tolerance and says
so. Every instance was certified in one arithmetic or the other, and a verdict
names which; \texttt{arithmetic\ =\ "exact"} refuses the fallback for a caller who
wants the stronger statement or nothing.

\subsection{The nineteen-solver portfolio}\label{the-nineteen-solver-portfolio}

Every solver here targets the objective of \eqref{eq:lap}, so the choice among
them is a performance question rather than a statistical one, and it depends on
the shape of the problem. The methods that reach it by scaling real costs onto
an integer grid are exact for the instance that scaling produces, which is the
supplied one where the costs are integers and a rounding of it otherwise;
\texttt{verify\_assignment()} is what decides whether a returned matching is optimal for
the matrix as supplied. We group them into four families and a
set of special-purpose routines.

\textbf{Augmenting-path methods} grow a matching one row at a time, each time finding
a shortest augmenting path in the residual graph under reduced costs. The
Hungarian method \citep{Kuhn1955, Munkres1957} is the classical member; the
Jonker-Volgenant algorithm \citep{JonkerVolgenant1987} adds column reduction and
auction-style initialisation, which is what makes it fast in practice on dense
problems despite an \(O(n^3)\) worst case. \texttt{lapmod} is the sparse variant, working
from an adjacency list so that forbidden pairs cost nothing to carry.

\textbf{Auction methods} \citep{Bertsekas1988} let rows bid for columns, raising prices
until no row wants to switch. They converge to optimality once the bid increment
\(\varepsilon\) falls below \(1/n\) for integer costs, and \(\varepsilon\)-scaling
gives the useful practical variants. Auction is naturally parallel and tolerant
of near-degenerate cost structures where augmenting-path methods spend time on
long paths.

\textbf{Cost-scaling methods} solve a sequence of increasingly precise problems,
each time halving the scale on which costs are rounded and reusing the previous
solution \citep{GoldbergKennedy1995}. The Gabow-Tarjan bit-scaling algorithm
\citep{GabowTarjan1989} belongs here. On a graph with \(|V|\) nodes and \(|E|\) edges it
reaches \(O(\sqrt{|V|}\, |E| \log(|V| C_{\max}))\) for integer costs, where
\(C_{\max} = \max_{i,j} |C_{ij}|\). A complete bipartite graph on \(n\) and \(m\) units
has \(|V| = n + m\) and \(|E| = nm\), so the bound there is
\(O(\sqrt{n + m}\; nm \log((n + m) C_{\max}))\).
We are not aware of a prior publicly available open-source implementation of
this 1989 bipartite assignment algorithm.

Bit-scaling runs on integers, so the conversion is part of the method and its
rule is stated rather than left to the caller. Finite costs are shifted so the
smallest is zero, multiplied by a scale factor and rounded, with the factor set
to keep the scaled costs clear of the sentinel marking a forbidden pair and the
path sums inside 64-bit arithmetic; a matrix whose range leaves no usable
factor is refused with the limit named. Integer costs convert exactly and the
optimum returned is the optimum of the instance as supplied. Costs that are not
integers are solved on the rounded instance, and the supplementary material
states the scale factor, the rounding rule and the bound on how far the
returned matching can sit above the optimum of the matrix as supplied.
Potentials are divided by the factor and shifted back, so they belong to the
original matrix, and whether the matching is optimal for it rather than only
for the rounded instance is a question \texttt{verify\_assignment()} answers.

\textbf{Network-flow methods} treat \eqref{eq:lap} as a minimum-cost flow problem on a
bipartite network \citep{Ahuja1993} and apply push-relabel \citep{GoldbergTarjan1988},
cycle cancelling, network simplex, or Orlin's strongly polynomial algorithm
\citep{Orlin1993}. These are the most general formulations, which makes them the
right base for variable-ratio designs where a group absorbs several units. On a
plain one-to-one problem they are the slowest choice.

\textbf{Special-purpose routines} cover cases where the general machinery is wasted:
exhaustive enumeration below nine units, Hopcroft-Karp bipartite matching
\citep{HopcroftKarp1973} when costs are binary or constant, Ramshaw-Tarjan
\citep{RamshawTarjan2012} for strongly rectangular problems where padding to a square
matrix would dominate the work, and bottleneck assignment, which minimises the
largest pair cost instead of the sum.

The binary case is worth stating in full, since a maximum-cardinality routine
carries no costs and the problem does. When every admissible edge holds the
same value, all perfect matchings cost the same and any maximum-cardinality
matching is optimal. When the finite costs are \(\{0, 1\}\), the routine runs
Hopcroft-Karp on the zero-cost edges alone: a perfect matching there totals
zero, which no assignment improves on. If that subgraph admits no perfect
matching the optimum needs at least one unit edge, and the problem passes to
the weighted solver on the original matrix. The specialised path answers only
where its answer is the optimum.

Two extensions return something other than a single optimal assignment. Murty's
ranking procedure \citep{Murty1968} with Lawler's partition scheme \citep{Lawler1972}
enumerates the k cheapest assignments, which lets a user see how much of the
total cost is forced and how much is arbitrary. Sinkhorn iteration \citep{Cuturi2013}
solves the entropy-regularised relaxation, returning a doubly stochastic
coupling rather than a permutation.

\subsection{Matrix-level interface and extensions}\label{matrix-level-interface-and-extensions}

The solver layer is reachable on its own, for a user who arrives with a cost
matrix. \texttt{lap\_solve()} and \texttt{assignment()} take a matrix and return the optimal
assignment, and \texttt{lap\_solve()} also accepts a long data frame of source, target
and cost columns, the natural shape when costs come from a database. The
\texttt{method} argument names any of the nineteen solvers and defaults to \texttt{"auto"};
it accepts twenty names, since \texttt{"ssp"} is a second spelling of \texttt{"sap"}.

\begin{verbatim}
set.seed(1)
cost <- matrix(sample.int(100, 25, replace = TRUE), nrow = 5)
res  <- lap_solve(cost)
res
\end{verbatim}

\begin{verbatim}
#> Assignment Result
#> =================
#> 
#> # A tibble: 5 x 3
#>   source target  cost
#>    <int>  <int> <int>
#> 1      1      4     7
#> 2      2      2    14
#> 3      3      1     1
#> 4      4      3    54
#> 5      5      5    44
#> 
#> Total cost: 120 
#> Method: bruteforce
\end{verbatim}

Dual potentials come from \texttt{assignment\_duals()}, which returns \(u\) and \(v\)
alongside the assignment, so passing that result to \texttt{verify\_assignment()}
certifies the matching without a second solve. Four extensions operate on the
same matrix: \texttt{lap\_solve\_kbest()} returns the k cheapest assignments in
increasing order of cost, \texttt{bottleneck\_assignment()} minimises the largest pair
cost rather than the sum, \texttt{lap\_solve\_batch()} solves a stack of independent
problems, optionally across threads, and \texttt{sinkhorn()} returns the entropy-regularised coupling, which \texttt{sinkhorn\_to\_assignment()} rounds to a
permutation.

\begin{verbatim}
kb <- lap_solve_kbest(cost, k = 3)
tapply(kb$total_cost, kb$rank, unique)
\end{verbatim}

\begin{verbatim}
#>   1   2   3 
#> 120 120 150
\end{verbatim}

The two cheapest assignments both cost 120 and the third costs 150, so the
optimum here is attained by two distinct pairings. A matching study reading
that output learns its matched sample is not unique, which stays invisible
when a solver returns one assignment. The solver portfolio and the
dispatcher that selects within it are measured across a grid of cost regimes,
admissibility patterns and aspect ratios in the empirical evaluation.

\section{Matching with couplr}\label{matching-with-couplr}

\subsection{A matching workflow}\label{a-matching-workflow}

The package ships \texttt{hospital\_staff}, a synthetic personnel dataset supporting a
treated/control workflow without external data: 200 units against 300, with
age, experience, hourly rate, department, certification level and a full-time
indicator.

\texttt{match\_couples()} is the front door. It takes the two data frames and the
covariates to match on and returns the optimal one-to-one pairing under the
requested distance; \texttt{auto\_scale\ =\ TRUE} screens the covariates for degenerate
variance and excessive missingness and picks a scaling method, so that
variables measured in years and in ordinal levels contribute comparably.

\begin{verbatim}
library(couplr)
data(hospital_staff)
treated <- transform(hospital_staff$nurses_extended,   id = nurse_id)
control <- transform(hospital_staff$controls_extended, id = nurse_id)
covars  <- c("age", "experience_years", "certification_level")

m <- match_couples(treated, control, vars = covars, auto_scale = TRUE)
m$pairs[1:3, c("left_id", "right_id", "distance", ".age_diff")]
\end{verbatim}

\begin{verbatim}
#> # A tibble: 3 x 4
#>   left_id right_id distance .age_diff
#>   <chr>   <chr>       <dbl>     <dbl>
#> 1 1       260         0.217        -1
#> 2 2       294         0.175         2
#> 3 3       212         0             0
\end{verbatim}

The call matches all 200 treated units at a total distance of
59.95. Beside the identifiers and each
pair's distance, \texttt{m\$pairs} carries a signed per-variable difference for every
matching variable, so which covariate drives a poor pair is visible without
recomputing anything; \texttt{m\$unmatched} holds what is left over on each side and
\texttt{m\$info} records the method, the pair count, the total distance and the
estimand.

\texttt{balance\_diagnostics()} computes standardised mean differences, variance
ratios and Kolmogorov-Smirnov statistics on the matched sample, and
\texttt{balance\_table()} formats them: experience and certification fall below the
conventional 0.1 threshold on the absolute standardised difference
\citep{Stuart2010} and age sits slightly above at 0.116. \texttt{join\_matched()} returns
the analysis-ready frame, one row per pair, with every column of both inputs
under \texttt{\_left} and \texttt{\_right} suffixes.

The matched frame is the input to an outcome analysis rather than the end of
one. A paired difference taken on it estimates an adjusted association, and
reading that as a causal effect is what rests on assumptions no matching method
can verify: conditional ignorability, positivity, and the stable unit treatment
value assumption. What matching can offer against the first is
\texttt{sensitivity\_analysis()}, which implements Rosenbaum bounds for the Wilcoxon
signed-rank statistic on one-to-one matched pairs \citep{Rosenbaum2002}, reporting
p-value bounds across a grid of the odds-ratio parameter \(\Gamma\) so that a
reader can judge how strong an unmeasured confounder would have to be to
overturn the result. The supplementary material carries both on the matched
sample above.

\subsection{Calipers, forbidden pairs and partial matching}\label{calipers-forbidden-pairs-and-partial-matching}

Three mechanisms restrict which pairs are admissible, and all three work by
writing non-finite entries into the cost matrix before it reaches the solver.
A global \texttt{max\_distance} forbids any pair beyond a distance ceiling. Per-variable
\texttt{calipers} forbid pairs whose raw difference on a named variable exceeds a
threshold, independently of the distance metric in use. Explicit forbidden pairs
are passed as non-finite entries by the user. Every solver recognises a
non-finite entry as a missing edge, so no returned pair ever violates a
constraint.

Tight constraints usually make \eqref{eq:lap} infeasible. couplr first drops
the rows and columns whose every edge is forbidden, reporting them unmatched.
The remaining submatrix can still fail Hall's condition, with several rows
competing for the same few columns and no assignment reaching every row.
Rather than return whichever pairs a heuristic happens to find, couplr
replaces the forbidden entries with a finite sentinel \(\sigma\) and re-solves
with the same optimal solver.

Write \([\ell, h]\) for the range of the admissible costs of the pruned
submatrix and \(k\) for the smaller of its two dimensions.

\textbf{Lemma 1.} \emph{If \(\sigma > k(|\ell| + |h|)\), then a padded assignment using more
sentinel edges than another costs strictly more.} The proof is in the
supplementary material.

A minimum-cost solve of the padded problem therefore uses as few sentinel
edges as any assignment can, which is a maximum-cardinality admissible
matching, and among the assignments attaining that number it minimises the
real cost. Dropping the sentinel pairs leaves the lexicographic optimum: the
largest admissible matching, cheapest among those of its size. couplr takes
\(\sigma = (k+1)(|\ell| + |h|) + 1\), which satisfies the hypothesis in every
sign case, and solves a maximisation by negation with the sentinel negated
alongside it. No sentinel enters a reported cost.

Above a padded objective of \(2^{53}\) couplr refuses the padding rather than
returning an ordering its arithmetic may not support: \texttt{match\_couples()} warns
and falls back to a greedy partial matching, saying that it is not optimal,
and \texttt{assignment(cardinality\ =\ "maximum")} errors and asks for rescaled costs
or an explicit \texttt{unmatched\_penalty}. That threshold is a precondition on the
solve and not a proof of it. What establishes that a particular padded solve
reached its optimum is the certificate, and the lemma carries that verdict
across to the lexicographic objective; the supplementary material sets out
both.

\begin{verbatim}
m_cal <- match_couples(treated, control, vars = covars, auto_scale = TRUE,
                       calipers = list(age = 3, experience_years = 2),
                       max_distance = 1.5)
\end{verbatim}

The caliper costs pairs: 197 of the 200 treated units keep a
partner. That is a caliper working, and the reason the unmatched identifiers
are returned rather than dropped, since the units that fail to match define
the population the estimate applies to. The constraints here are severe,
leaving 2,173 admissible pairs out of 60,000, so the padded solve is what
produces this result; a greedy pass over the same edges matches 180 units
rather than 197.

Blocking is the other structural constraint. \texttt{block\_id} restricts matching to
levels of a factor, and \texttt{matchmaker()} builds blocks when no grouping variable
exists, from an existing variable or by k-means or hierarchical clustering on
chosen block variables. The solver runs once per block, and blocked designs
report per-block balance so imbalance inside one stratum is not averaged away
across the others.

\subsection{Matching designs}\label{matching-designs}

Beyond one-to-one matching, \texttt{match\_couples()} takes \texttt{ratio} for k-to-one designs
and \texttt{replace} for matching with replacement. Five further designs have their own
entry points, each returning an object in the same family. \texttt{ps\_match()} estimates
a propensity score \citep{RosenbaumRubin1983} and matches on it, with the caliper
expressed in standard deviations of the linear predictor. \texttt{full\_match()} assigns every unit to a matched group of variable size, so no
unit is discarded, solving a minimum-cost flow problem with Johnson potentials
\citep{Johnson1977} on its optimal path and a greedy two-pass heuristic on large
problems. \texttt{cem\_match()} implements coarsened exact matching
\citep{IacusKingPorro2012}, and \texttt{subclass\_match()} divides the sample into
propensity-score subclasses, returning subclass weights for the average
treatment effect on the treated or on the whole sample.

\texttt{cardinality\_match()} targets the cardinality-matching objective
\citep{Zubizarreta2014}, the largest matched sample meeting prespecified balance
constraints, with total distance ordered after cardinality; which solver answers it follows
from the kind of balance asked for. Fine and refined covariate balance
\citep{RosenbaumRossSilber2007, Pimentel2015}, stated through \texttt{fine} and
\texttt{refined}, are representable in the network as category nodes, so such a call
reaches a single minimum-cost flow solve returning the largest balanced sample
with a dual certificate at polynomial cost. Linear moment constraints, a
finite \texttt{max\_std\_diff} or an explicit \texttt{moments} argument, are not, and are
dualised and searched by branch and bound under \texttt{node\_limit} and \texttt{time\_limit}.
Either way the result reports \texttt{n\_matched} beside \texttt{best\_possible}, the largest
sample the constraints admit, the \texttt{gap} between them, whether the answer is
\texttt{certified}, and what the search \texttt{stopped\_on}. A pruning loop is available as
\texttt{engine\ =\ "heuristic"} and reports its gap as unknown, which is the honest
reading of what a heuristic establishes.

\subsection{Interoperability}\label{interoperability}

couplr fits into a workflow built around the established packages.
\texttt{as\_matchit()} converts a couplr result into a \texttt{matchit} object and
\texttt{match\_data()} returns data in the layout \CRANpkg{MatchIt} users expect, so
downstream code keeps working; \texttt{bal.tab()} methods are registered for the
couplr classes, so \CRANpkg{cobalt} produces its usual balance output; and
weighted matched data flows into \CRANpkg{marginaleffects} \citep{marginaleffects}.

\section{Common optimisation architecture}\label{common-optimisation-architecture}

Three of the package's capabilities are consequences of one decision.
Compiling every flow-representable design into a single flow model is what
makes node potentials available from all of them; the potentials are what
certify a solution; a certificate is what makes it safe to optimise over a
restricted set of pairs; and a certified restricted solve is what a warm start
resumes from:

\[\text{common flow duals} \rightarrow \text{certificates}
\rightarrow \text{safe edge generation} \rightarrow \text{warm-started paths}.\]

\subsection{Flow compilation}\label{flow-compilation}

Underneath the designs of the previous section there is one object. It has
nodes, each carrying a signed supply, and arcs, each carrying a cost and a
capacity range. As Table \ref{tab:designs} sets out, a one-to-one match, a
k-to-one match, matching with replacement and a full match differ in the bounds
their arcs receive, not in the code that solves them. Formulating full matching as a minimum-cost flow problem is established
\citep{HansenKlopfer2006}; what the flow model contributes here is that every
design in the package that a flow can represent compiles into it, so one
solver and one certificate serve all of them, with the warm-start machinery
shared where the public interface exposes it. Those are the
six rows of the table: one-to-one, k-to-one, with replacement, blocked and
exact, full matching, and fine and refined balance. The designs that sit
outside it sit outside because their answer is not a flow: \texttt{cem\_match()}
coarsens and strata are read off directly, \texttt{subclass\_match()} cuts the
propensity score into subclasses, and the moment constraints of
\texttt{cardinality\_match()} are dualised and searched by branch and bound rather
than represented in a network.

\begin{longtable}[]{@{}ll@{}}
\caption{\label{tab:designs} How each matching design is stated in the internal flow
model. The designs differ in the capacity bounds their arcs receive and in how
many networks are built, rather than in which solver runs.}\tabularnewline
\toprule\noalign{}
Design & Flow formulation \\
\midrule\noalign{}
\endfirsthead
\toprule\noalign{}
Design & Flow formulation \\
\midrule\noalign{}
\endhead
\bottomrule\noalign{}
\endlastfoot
One-to-one matching & Unit-capacity bipartite flow \\
k-to-one matching & Capacitated bipartite flow, treated supply k \\
Matching with replacement & Control capacities above one \\
Blocked and exact matching & Independent networks per stratum \\
Full matching & Lower and upper capacitated circulation \\
Fine and refined balance & Category nodes with capacity bounds \\
\end{longtable}

Compiling a design, solving it, and checking the answer stay three separate
steps, because the third is worth something only while it is separate: a
solver status is not itself an independently checkable certificate. The flow and the potentials
go into \texttt{verify\_flow()} as inputs to a check, and the arcs the flow is checked
against are stated independently of it.

\begin{verbatim}
prob <- list(n_nodes = 4, supply = c(2, 1, -2, -1),
             arcs = data.frame(tail = c(1, 1, 2, 2), head = c(3, 4, 3, 4),
                               lower = 0, upper = 2, cost = c(1, 3, 2, 1)))
cert <- verify_flow(c(2, 0, 0, 1), prob, potential = c(0, 0, 1, 1))
c(certified = cert$certified_optimal, gap = cert$duality_gap)
\end{verbatim}

\begin{verbatim}
#> certified       gap 
#>         1         0
\end{verbatim}

Two solvers read that network. One is successive shortest paths with Johnson
potentials \citep{Johnson1977}, reachable as \texttt{method\ =\ "csflow"}. The other, \texttt{method\ =\ "push\_relabel"}, is cost-scaling push-relabel in the
successive approximation of \citet{GoldbergTarjan1990}: a sequence of
\(\varepsilon\)-optimal flows, each phase dividing \(\varepsilon\), saturating the
arcs the smaller \(\varepsilon\) no longer admits, and clearing the excess with
two local operations. A push moves excess along an arc of negative reduced
cost; a relabel lowers the price of a node holding excess and having no such
arc, by exactly the amount that gives it one. Below \(\varepsilon = 1/(n+1)\) an
\(\varepsilon\)-optimal flow on integer costs is optimal, which ends the
scaling, so real costs are scaled to integers first.

Tolerance is where a flow certificate can overstate what it knows, and the
scale of the potentials sets it. A design that ranks objectives lexicographically
by weighting them, as \texttt{cardinality\_match()} ranks cardinality above balance
above distance, reaches potentials in the millions, where one unit in the last
place is around \(10^{-9}\) and an exactly optimal flow cannot meet an absolute
tolerance of that size. The comparisons therefore scale with the arc's own
magnitude, and the widest tolerance any comparison used comes back as
\texttt{dual\_tolerance}, so the verdict names the resolution it was reached at.

Node potentials therefore come back from every design rather than from the
one-to-one case alone. Those potentials are what the next section prices
against.

\subsection{Certified implicit edge generation}\label{certified-implicit-edge-generation}

A matching problem on \(n\) treated units and \(m\) controls has \(nm\) admissible
pairs, and materialising them is what sets the size ceiling: at 50,000 against
100,000 the dense cost matrix alone is 40 GB. The lazy representation removes
the matrix while leaving the search over all \(m\) columns intact. What
\texttt{memory\_mode\ =\ "implicit"} changes is the problem the solver ranges over: it
optimises over a restricted pair set and certifies the result against the
complete one, so the pairs outside that set are priced rather than searched.

The loop is column generation \citep{LubbeckeDesrosiers2005} over the pair variables
of \eqref{eq:lap}. A restricted master problem holds a subset of them, its
optimal duals price the rest at \(C_{ij} - u_i - v_j\), and the restricted
solution solves the full problem as soon as no omitted variable prices negative.
One thing separates it from the textbook setting: the omitted variables are
enumerated rather than returned by an optimisation oracle, because a pair
variable is an index rather than a combinatorial object to be searched for.
Write
\(E_t \subseteq \{1, \dots, n\} \times \{1, \dots, m\}\) for the candidate set at
round \(t\) and \(\epsilon\) for the pricing threshold.

\begin{enumerate}
\def\labelenumi{\arabic{enumi}.}
\tightlist
\item
  Seed \(E_0\) with each row's \(w\) nearest admissible columns, \(w\) read off \(m\)
  as \(6\lceil \log_2 m\rceil\) and capped at \(m\).
\item
  Solve \eqref{eq:lap} restricted to \(E_t\) with the flow model of the previous
  section, warm-started from the incumbent flow and prices when \(t > 0\). If it
  comes back short of a complete matching, repair the shortfall as below and
  repeat this step.
\item
  Read the master's potentials as assignment duals \((u, v)\).
\item
  Price the omitted pairs: \(\bar{C}_{ij} = C_{ij} - u_i - v_j\) for admissible
  \((i,j) \notin E_t\).
\item
  Stop if \(\bar{C}_{ij} \ge -\epsilon\) everywhere.
\item
  Otherwise set \(E_{t+1} = E_t \cup V_t\), where \(V_t\) holds each row's most
  negative violators, and return to step 2.
\end{enumerate}

The stopping rule is the certificate.

\textbf{Proposition 1.} \emph{Let \(X\) solve \eqref{eq:lap} restricted to a pair set \(E\)
with optimal duals \((u, v)\), and suppose \(C_{ij} - u_i - v_j \ge 0\) for every
admissible \((i,j) \notin E\). Then \(X\) is optimal for the complete problem.}

\emph{Proof.} Optimality on \(E\) gives complementary slackness for \(X\) and \((u,v)\)
there, so every matched pair is tight and every unmatched column has \(v_j = 0\),
and the cost of \(X\) is \(\sum_i u_i + \sum_j v_j\). The hypothesis extends dual
feasibility from \(E\) to every admissible pair, so \((u, v)\) is feasible for
\eqref{eq:duals} on the complete problem and that sum is a lower bound on the
cost of every assignment of it. \(X\) is one of them and attains the bound.
\(\square\)

Nothing in the argument is numerical: the hypothesis is the sign of
\(\bar{C}_{ij}\) on the omitted pairs, and the conclusion is equality, not
proximity. In practice four quantities weaken it, and the supplementary
material accounts for them separately: the gap \(\delta\) the master's own
certificate reports, the pricing threshold \(\epsilon\), the sign condition on
the column potentials, and the allowance the pricing bound carries. Their sum
bounds the matching's excess over the complete optimum. At
\(\epsilon = 10^{-9}\) and the largest shape in the benchmark below, the
threshold's contribution \(n\epsilon\) is \(2 \times 10^{-5}\) in the cost
units, and every run reported there came back with \(\delta = 0\).

The loop terminates because each round adds at least one pair \(E_t\) did not
hold, so \(|E_t|\) strictly increases under the bound \(nm\). That bound is the
worst case: duals that price most of the grid in make the loop build the
complete graph and pay for the rounds that got there. \texttt{max\_rounds}, 60 by
default, is a guard rather than the termination argument, and a run reaching it
returns \texttt{iteration\_limit} rather than a certificate.

A pricing sweep stores nothing: the reduced cost of a pair is consumed on the
spot, so the memory the mode holds is the candidate set, \(nw\) pairs from the
seed plus what the rounds add. How it reads the pairs depends on the cost
source. Where the metric admits a ball bound the columns are held in a metric
tree, and a subtree whose bound cannot beat the row's current threshold is
discarded without being visited; a source with no such bound is read one pair
at a time, at \(O(nm)\) arithmetic to the sweep. Mahalanobis distance is
Euclidean after whitening and always pays for the bound, being quadratic in the
covariates, while the metrics linear in them take a tree only in low dimension.
Where the bound prunes, the prune certifies that its subtree holds no violating
edge, and the prunes with the leaves examined certify the pricing result. The
supplementary material states the bound, the summaries a node carries and the
allowance every bound is lowered by. Pruning does not put a round below one
pass over the pairs, since the seed evaluates what it seeds with and a pair
priced in one round is priced again in the next, so \texttt{edges\_evaluated} can
exceed \texttt{possible\_edges} where the tree is working; what the loop saves is the
graph.

Rows contribute at most \texttt{keep\_per\_row} violators each, five by default, so a
round adds at most \(5n\) pairs and never fewer than one. Selection is by strict
improvement on a row's worst kept reduced cost over a scan running columns
ascending, so ties keep the lower column index and the same input gives the
same candidate set at every round.

A master that comes back short of a complete matching has no dual solution to
price with, so the shortfall is repaired rather than priced. Hall's condition
\citep{Hall1935} names a row set \(S\) whose neighbourhood \(N(S)\) in the restricted
graph is too small to go round. Every restricted edge out of \(S\) already lands
in \(N(S)\), so only columns outside it can move the deficiency, and the re-seed
reads the full source over the rows of \(S\) and takes columns from outside
\(N(S)\). When there are none, \(N(S)\) is that neighbourhood in the complete problem too,
and the result is infeasible with \((S, N(S))\) as its witness, naming the units
short of partners rather than only the fact.

So the sparse solution is optimal for a problem that was never built, and
\texttt{certify\ =\ TRUE} runs that check and returns it.

Two quantities in \texttt{\$search} describe what that cost. \texttt{candidate\_edges} against
\texttt{possible\_edges} is the graph built against the complete one that was not, and
\texttt{edges\_evaluated} counts distance evaluations over all rounds.

Round count is the quantity to control, since each round prices every omitted
pair again, and it is what the seed rule of step 1 is sized for: a seed too
narrow buys what it needs a round at a time, a seed too wide pays for columns
the matching never uses. The rule clears the two-round floor, one seed and one
sweep that finds nothing to add, on the problems below, and a run reports the
width it used as \texttt{\$search\$seed\_width}.

\begin{verbatim}
m_imp <- match_couples(treated, control, vars = covars, auto_scale = TRUE,
                       memory_mode = "implicit", certify = TRUE)
m_imp$certificate$certified_optimal
\end{verbatim}

\begin{verbatim}
#> [1] TRUE
\end{verbatim}

\begin{verbatim}
unlist(m_imp$search[c("seed_width", "n_rounds",
                      "candidate_edges", "possible_edges")])
\end{verbatim}

\begin{verbatim}
#>      seed_width        n_rounds candidate_edges  possible_edges 
#>              54               2           10800           60000
\end{verbatim}

The certificate the loop returns is the numerical one, which is why the
\(\epsilon\) version of Proposition 1 is the statement it carries. Dual
feasibility over the pairs the master omits comes from the pricing bounds, which
ask whether anything prices below \(-\epsilon\) rather than evaluating each
omitted pair, so the exact conclusion of the certification section is not
reachable from a scan assembled that way, and the certificate reports that as
its arithmetic.

Constraints tighten the loop rather than complicating it, since a caliper or a
distance ceiling is checked before a distance is computed and removes pairs from
the set to be priced.

Two neighbours are worth separating from this. \citet{Abeywickrama2021} also decline
to build a complete cost matrix, deferring exact edge costs behind an
inexpensive, application-specific lower bound, so what they save is
arithmetic. The bound here is geometric rather than application-specific, and
what the loop saves first is the graph it never builds; where the metric
admits a ball bound the pricing also declines to evaluate a pair, which is
where the two meet. Graded matching \citep{YuRosenbaum2022} also enlarges a sparse network
only as far as it has to, but a grade is fixed before the match, while the
reduced cost that admits a pair here is a property of the solution the master
has just returned, so the loop stops on a statement about the complete graph
rather than on the grades running out.

The mode is available on \texttt{match\_couples()} and \texttt{assignment()}, and it is
requested rather than selected: \texttt{memory\_mode\ =\ "auto"} does not choose it,
because the worst case above is a real one. Its scope is the unit-capacity assignment. A full match's column nodes carry
capacities above one, so a column dual there is not the assignment dual this
pricer reads. A general flow arc does carry a reduced cost of its own, from
the potentials at the nodes it joins, so pricing one is available in
principle; it is this implementation that stops at the assignment. How small a graph the loop
ends up holding, how many pricing rounds it takes and how it behaves on
geometries built to be hostile to it are measured in the empirical
evaluation.

\subsection{Warm-started matching paths}\label{warm-started-matching-paths}

A caliper is rarely known in advance. The usual practice is to sweep it, watch
what the matched sample and its balance do, and choose a point, which means
solving the same problem many times over. \texttt{match\_path()} solves the sweep as
one sequence.

\begin{verbatim}
p <- match_path(treated, control, vars = covars, auto_scale = TRUE,
                vary = "max_distance", values = c(1.0, 1.2, 1.5, 2.0, 3.0))
p$path[, c("max_distance", "n_matched", "total_distance", "certified")]
\end{verbatim}

\begin{verbatim}
#> # A tibble: 5 x 4
#>   max_distance n_matched total_distance certified
#>          <dbl>     <int>          <dbl> <lgl>    
#> 1          1         200           62.6 TRUE     
#> 2          1.2       200           61.5 TRUE     
#> 3          1.5       200           59.9 TRUE     
#> 4          2         200           59.9 TRUE     
#> 5          3         200           59.9 TRUE
\end{verbatim}

The sweep shows where the constraint stops binding: the total distance falls
until 1.5 and does not move above it, since a cut that wide forbids nothing
the unconstrained optimum uses. Reading that off a path is the point of
solving the sweep rather than one value.

Each point resumes from the matching the point before it found. Raising the
distance cut admits pairs while leaving that matching feasible, so what
remains is to repair reduced-cost violations on the arcs that just became
admissible, which is one round of the pricing loop of the previous section; an
independent call at that value pays for a solve from cold. One mechanism
serves the sweep and the single solve. A value whose admissible pairs cannot
support a complete matching is reported rather than approximated, coming back
with no matching and the Hall witness beside it, so a sweep that starts below
the feasibility threshold says where that threshold is. The ascending
direction is a correctness requirement rather than a convention: a descending
sweep withdraws pairs the current matching may be standing on, which breaks
feasibility and needs a repair phase the loop does not have, so a descending
sequence is refused with a message saying why.

The return value is a \texttt{couplr\_path}, whose \texttt{\$path} carries a row per point
with the matched count, total distance, seconds, rounds and whether it was
certified, and whose \texttt{\$balance} carries a row per point per variable. \texttt{certify\ =\ TRUE} is the default here, and a point's certificate is what says its
matching is the optimal one at that value rather than the one the warm start
happened to reach. What a sweep costs against solving each point
independently is measured in the empirical evaluation.

\subsection{Memory representations}\label{memory-representations}

The dense representation of the cost matrix is the binding constraint on problem
size long before runtime is. A 100,000 by 100,000 problem needs 80 GB as doubles;
the same units described by 100 covariates need 80 MB. couplr exposes this as \texttt{memory\_mode}, with values \texttt{"dense"}, \texttt{"lazy"},
\texttt{"implicit"} and \texttt{"auto"}, differing in what the solver ranges over. The dense
path holds the complete matrix; the lazy path holds none and still considers
every pair, evaluating each distance when it is asked for, over Jonker-Volgenant and auction with the built-in metrics; the implicit path runs the
edge-generation loop described above. All three solve the same problem, and on
the instances reported below they returned the same pairing wherever the pairing
was compared.

Under \texttt{"auto"} the package decides. Below ten million cells, about 80 MB
dense, it skips the check, since probing free memory on every ordinary problem
would be pure overhead. Above that it estimates what a dense solve peaks at
rather than what the matrix occupies, at 10 times the raw cell
bytes, and compares that against available system memory read per platform.
The empirical evaluation reports the peaks that multiplier sits above. If the
estimate exceeds half of what is available, \texttt{"auto"} switches to the lazy path
where one exists and warns where
one does not, naming blocking, greedy matching or a smaller problem. Detection
failure is treated as unknown rather than as success: a fixed 4 GB threshold
takes over, so the guard never silently disappears on an unrecognised
platform.

\section{Implementation and verification}\label{implementation-and-verification}

The package is organised in three layers. The first turns user input into a
cost source: it validates the data frames, screens covariates, applies scaling
and weights, and settles how distances are obtained and which pairs are
admissible. The second compiles the requested design into the flow model and
solves it, which for the plain one-to-one case is \eqref{eq:lap}. The third
builds result objects and their methods. The object model, the C++
organisation of the nineteen solvers and the dependency footprint are in the
supplementary material.

\subsection{Automatic solver selection}\label{automatic-solver-selection}

\texttt{method\ =\ "auto"} inspects the problem and picks a solver, and the rules are
few.

\begin{enumerate}
\def\labelenumi{\arabic{enumi}.}
\tightlist
\item
  Problems with at most eight rows and eight columns go to exhaustive
  enumeration, which is exact and faster than setting up a general solver.
\item
  Cost matrices whose finite entries are constant or lie in \(\{0, 1\}\) go to
  the Hopcroft-Karp style routine, which exploits the absence of a real cost
  scale.
\item
  Everything else goes to Jonker-Volgenant.
\end{enumerate}

Sparsity and aspect ratio carried rules of their own, sending a mostly
forbidden matrix to \texttt{lapmod} and a wide one to shortest augmenting paths. Both
were slower than the default in the regime they were written for, so both were
removed; Table \ref{tab:regimes} is the measurement, and both solvers stay
reachable by name.

Rule 2 reads the entries, and so does the rejection of \texttt{NaN} inputs that runs
for every method, so choosing a solver costs one C++ pass over the cost matrix,
returning every quantity the rules need without allocating and without a second
scan. That leaves the inspection a linear read in front of a superlinear
solve, so its share of the runtime falls as problems grow, and the dashed line in
Figure \ref{fig:solver-bench} sits on the solver the dispatcher selects. The
pass is skippable for callers who already know the shape of their problem, by
naming the method.

\subsection{Independent verification}\label{independent-verification}

Certificates are the check that survives leaving the test suite. \texttt{verify\_assignment()} and \texttt{verify\_flow()} run in \(O(nm)\) and
\(O(\text{arcs})\) respectively without solving anything, so a user can run them
on their own problem, in exact arithmetic or at their own tolerance, against a
solution the package produced.

\subsection{Testing and reproducibility}\label{testing-and-reproducibility}

A smoke test establishes that a function returns an object of the right shape;
for a solver the standard is whether the assignment is the optimal one. Every
optimal solver is checked against exhaustive enumeration on randomised
instances covering integer and fractional costs, square and rectangular
shapes, minimisation and maximisation, and forbidden edges. The reference is
the same \texttt{bruteforce} solver the dispatcher uses below nine units, so the
verified and the shipped path are one implementation. Constrained matching is
checked against an enumeration over all partial matchings of a small instance,
which supplies both halves of its lexicographic objective. The
paths that avoid materialising the problem are checked against the path that
materialises it: the implicit loop and the dense solve must return the same
total wherever the dense solve runs at all, and a warm-started point must return
what an independent solve at that value returns. Pair identity is the stronger
check and is reported where it holds rather than required, since an instance
with several optima admits more than one answer of equal cost; what correctness
asks for is a certified objective the two agree on. The next section reports
both comparisons on the sizes it covers.

The suite covers the R layer, the solver bodies compiled outside R, the C++
wrappers, dispatch rules and the matching designs.

Every number the next section reports comes from a script that ships with the
article. Each caches its CSV under \texttt{paper/} and re-measures only when that file
is moved aside; \texttt{paper/run\_bench\_suite.sh} runs them in order on one machine,
which is how these numbers were produced, and the scaling, implicit,
implicit-grid, regime and memory scripts take \texttt{-\/-quick} for a reduced grid. The
environment block beside the CSVs records the machine, the R version, the
package version and the commit.

\section{Empirical evaluation}\label{empirical-evaluation}

\subsection{Edge generation, robustness and memory}\label{edge-generation-robustness-and-memory}

Table \ref{tab:implicit-table} reports the three memory modes on those same
problems. The dense and lazy arms are pinned to Jonker-Volgenant, so they differ
only in representation; the implicit arm carries the flow model as its
restricted master by construction, which is a difference between the arms and is
the definition of the mode rather than a choice.

\begin{table}

\caption{\label{tab:implicit-table}One-to-one optimal Mahalanobis matching by memory mode, wall-clock seconds on a single core of an Apple M4 Pro. Graph built is the arcs the implicit loop ended up holding as a share of the complete problem's. Distances is its evaluations over all rounds as a multiple of the complete pair count, so a pair priced twice counts twice. All three arms were timed in one session; the dense arm was run at the four sizes where its pairing can be compared against the loop's, and the scaling table below carries dense timings at the two largest sizes. Every cell returned the same total distance.}
\centering
\begin{tabular}[t]{lrrrrrr}
\toprule
Problem size & dense & lazy & implicit & Graph built & Distances & Rounds\\
\midrule
167 + 333 & 22 ms & 8 ms & 12 ms & 16.22\% & 2.00 & 2\\
667 + 1,333 & 161 ms & 53 ms & 56 ms & 4.95\% & 2.00 & 2\\
1,667 + 3,333 & 837 ms & 351 ms & 299 ms & 2.16\% & 1.96 & 2\\
3,333 + 6,667 & 3.14 s & 1.59 s & 1.09 s & 1.17\% & 1.88 & 2\\
6,667 + 13,333 & not run & 6.38 s & 3.9 s & 0.63\% & 1.71 & 2\\
\addlinespace
16,667 + 33,333 & not run & 62 s & 20.2 s & 0.29\% & 1.42 & 2\\
\bottomrule
\end{tabular}
\end{table}

Three things in that table go together. The loop holds a small graph, 16.22\% of the pairs at the smallest size and 0.29\% at the largest, a
share that falls as the problem grows since the pairs a matching can use grow
more slowly than the pairs a problem has. It pays for that in arithmetic, at
1.42 to 2.00 times the complete pair count in distance
evaluations, so at eight covariates it computes more distances than one full
sweep would. And it is still the fastest arm at four of the six sizes, because
what it avoids is not the arithmetic but the graph. At
16,667 + 33,333 it finishes in
20 seconds against
62 for the lazy path. At the two
smallest sizes the lazy path is faster; at 667 + 1,333 it takes
53 ms against
56 ms, which is where the
loop's fixed costs are still visible against a problem this size.

Every run in the table settled in 2 rounds, one seed and one sweep
that found nothing to add and certified. Every one came back certified with a
duality gap of zero. On the four sizes where the dense solve can also be run,
the two agree to within 3.4e-13, and the
pairing is identical unit by unit, so the loop returns the assignment rather
than an assignment of the same cost.

The implicit column is a representation and a solver at once, so comparing it
with the dense column reads both differences together. At 3,333 + 6,667 the
materialised matrix takes 3.14 seconds
under Jonker-Volgenant and 4.17
under the flow solver the restricted master belongs to, against 1.09 for the loop, so what the loop
saves is not what a change of solver would give. The supplementary material
carries the decomposition.

One cloud is not a characterisation, and the table above reports one. The
supplementary material sweeps the loop around that cloud one factor at a time:
covariate counts from 2 to 32, seven point
clouds including clustered points, a displaced treated group, coordinates on a
coarse lattice so that distances tie, one contested core where every treated
unit's nearest controls are the same few, and a shell where every
treated-to-control distance is nearly equal, three caliper settings including
the tightest value each problem is still feasible under, and five sizes, each at
several seeds. Across the 20 cells that grid holds, the loop settled in
2 to 10 rounds and ended holding between
0.29 and 73.32 percent of
the complete pair count, at 0.24 to
6.67 times that count in distance evaluations.
All 20 cells certified, at a largest absolute duality gap of 3.64e-12, and 18 of the 18 cells small enough to run the dense
path beside the loop returned the same pairing unit by unit. The margin over the lazy path is where the clouds separate: it runs from
11.13x on shifted to
1.09x on heavy\_tailed.

Two rounds is what an ordinary cloud costs and it is not what the loop is
bounded by, so the grid's hostile cell is worth reading on its own. In the
contested core, where every treated unit's nearest controls are the same few,
the duals travel a long way before they stop pricing pairs in. At 10,000 units it takes 9
to 10 rounds, ends holding 73 percent of the complete pair count rather than a
fraction of one percent, and evaluates 6.7 times that count in distances. It still certifies,
and it is still 3.2 times faster than the
lazy path on this cloud. What it shows is the worst case of the previous
section happening: the graph the loop ends up holding is most of the one it
set out to avoid, and it has paid for the rounds that built it. This is why
the mode is requested rather than selected.

\FloatBarrier

The claim that memory rather than time is what bounds a dense matching has so
far been made with edge shares, which are a proxy for it. Peak resident set size
is the quantity itself, and it is a property of a process rather than of an
expression, so each arm here is a fresh R session that loads what it needs, runs
one matching and exits, measured from outside. A session that loaded the same
packages and matched nothing is the floor each arm is read against.

At 6,667 + 13,333 the dense path peaks 6234 MB
above that floor, the lazy path 22 MB and the
loop 128 MB, against
27428 MB for \CRANpkg{optmatch} and
8640 MB for \CRANpkg{MatchIt}. Two readings follow. The loop is not the smallest of the three: the lazy path
holds 22 MB against the loop's 128 MB, since it stores no graph at all
while the loop stores the candidate set it has built. What the loop buys over
it here is time, and the ceiling both of them remove is the dense path's. The
second reading is the guard's. Building the matrix and stopping there accounts
for 1095 MB on a matrix whose entries
occupy 711 MB, so holding it costs 1.5 times its cells and solving over it costs 8.8 times them. That gap is preparation copies
and the solver's own workspace, and it is what the automatic mode has to read:
a guard sized on the matrix would let a dense solve start on a machine it does
not fit. \texttt{estimate\_dense\_matrix\_mb()} estimates the matrix at
4 times the raw cell bytes and \texttt{estimate\_dense\_solve\_mb()} the
peak of the solve at 10 times them, and \texttt{memory\_mode\ =\ "auto"}
compares the second against available memory. Neither multiplier is fitted to
one cell. Across 5,000, 10,000 and 20,000 units the solve peaked at
9.5, 7.2 and 8.8 times the cell bytes and the matrix at 1.8, 1.7 and 1.5
times them, so each default sits above the largest of those rather than at any
one of them. That is the direction the guard has to err in, since it exists to
refuse a solve that will not fit rather than to predict where the peak lands.
Every arm and size is in the supplementary material.

\FloatBarrier

\subsection{Warm-started paths}\label{warm-started-paths}

Table \ref{tab:path-table} reports a sweep of 20 caliper values
solved as one path against the same 20 values solved
independently. The grid is the data's own: its tight end is the tightest caliper
the problem still admits a complete matching under, found by bisection, and its
wide end is the median treated-to-control distance, which admits half the pairs
and constrains nothing.

The table leads with the wall clock for the whole sweep, which is the number a
user meets: it carries the feasibility check, the balance reading each point
takes, the conversions in and out of the solver and the rest of the R-level
work. Both sides pay that per point, since an independent solve here is a path
of one value rather than a separate entry point, so the comparison is like for
like. The solver's own time sits beside it as a diagnostic, saying how much of
the saving is the warm start rather than the sweep's shared setup.

\begin{table}

\caption{\label{tab:path-table}A caliper sweep solved as one warm-started path against the same values solved independently, end-to-end wall clock for the whole sweep on a single core of an Apple M4 Pro. Cold/path is the independent sweep's time over the warm-started path's, so a value above one is the path running faster. Solve-time ratio is the same comparison on the seconds the solver reports for itself, summed over the points, which excludes the R-level work both sides do. Rounds is the pricing rounds each side took, summed over the sweep.}
\centering
\begin{tabular}[t]{lrrrrr}
\toprule
Problem size & As a path & Independently & Cold/path & Solve-time ratio & Rounds\\
\midrule
167 + 333 & 0.057 s & 0.136 s & 2.39x & 2.88x & 22 / 40\\
667 + 1,333 & 0.489 s & 1.04 s & 2.12x & 2.15x & 23 / 40\\
1,667 + 3,333 & 2.88 s & 6.12 s & 2.13x & 2.13x & 21 / 40\\
6,667 + 13,333 & 40.9 s & 84.6 s & 2.07x & 2.07x & 21 / 40\\
\bottomrule
\end{tabular}
\end{table}

The saving comes from the rounds column. A warm-started point starts from a
matching that is still feasible at the new value, so what it has to repair are
the reduced-cost violations on the arcs the wider cut just admitted, which the
loop reaches in about one round per point. A point solved independently starts
from a seed and pays for the seeding sweep as well.

Correctness is the part of this comparison that has to hold before the timing
means anything. Every point on the path returns the status and the matched count
its independent solve returns, and the totals agree to
0 relative, so the path is the same
sequence of matchings that solving each value from cold would produce.

\FloatBarrier

\subsection{Solver portfolio and automatic dispatch}\label{solver-portfolio-and-automatic-dispatch}

Figure \ref{fig:solver-bench} reports median wall-clock solve time against
problem size for all nineteen solvers, grouped by family. We draw costs as
integers uniformly from \([1, 10{,}000]\) on square matrices and report medians of
five replicates on a single core of an Apple M4 Pro. The slower families are
capped at smaller sizes so that the scalable solvers can be measured to
\(n = 5000\) within a reasonable budget. All timings in this section were measured on 2026-08-30 under R 4.5.3 with couplr 1.7.0 at commit a63b760,
\CRANpkg{optmatch} 0.10.8 and \CRANpkg{MatchIt} 4.7.2.

\begin{figure}

{\centering \includegraphics[width=1\linewidth,alt={Five panels of log-log line plots showing solve time in milliseconds against problem size n from 4 to 5000. Solve time rises with problem size in every panel. The Jonker-Volgenant panel reaches n equals 5000 in roughly one second, the fastest of the families. Auction, cost-scaling and flow-based solvers are one to three orders of magnitude slower at matched sizes and stop at n equals 1000 or 2000. The dashed dispatcher line lies on top of the fastest solver in each panel.}]{rjournal_files/figure-latex/solver-bench-1} 

}

\caption{Median wall-clock solve time against problem size for the nineteen assignment solvers in couplr, grouped by algorithm family on shared log-log axes. The two special-purpose solvers in the Other panel run on their own inputs and are not comparable to the rest: HK-01 is timed on binary cost matrices, and Brute-F only up to n = 8. The dashed grey line repeated in every panel is the automatic dispatcher on the uniform integer costs.}\label{fig:solver-bench}
\end{figure}

Jonker-Volgenant with its warm start is the fastest general-purpose solver at
every measured size, which is why it is the dispatcher's default. The auction and cost-scaling families are competitive at small sizes and fall
behind on dense uniform costs, where their advantage, tolerance of degenerate
structure and a better asymptotic bound for Gabow-Tarjan, does not pay.
General flow formulations are the slowest, as expected of machinery solving a
larger problem than the one posed. The two special-purpose routines run on the inputs they exist for, exhaustive
enumeration below nine units and the Hopcroft-Karp style solver on binary
costs, so their timings are not comparable with the rest and are shown only to
establish that each is practical on its own input. These orderings hold for
dense uniform integer costs on square matrices, which is the whole of what the
figure establishes, and they fix the dispatcher's default.

The grid behind Table \ref{tab:regimes} is what reduced the rules to three.
It crosses the properties the solver figure holds fixed: whether the finite
costs have a scale worth exploiting, how much of the matrix is forbidden and
whether what remains is one component, and how far from square the problem is.
Eight cost regimes, from uniform integers through binary, constant, heavily
tied and heavy-tailed costs to Euclidean distances over uniform and clustered
points, are crossed with six admissibility patterns, complete, unstructured at
four densities from 60 percent of the entries finite down to 1 percent, and
four disconnected components, and with three aspect ratios from square to one
row per ten columns. That crossing is the base tier. The large tier repeats a
subset of it at three times the size rather than the whole: the regimes and
patterns where a rule that holds at 500 rows might fail at 1500, against the
solvers that reach that size inside the budget. The table therefore covers 189 cells, 144 of them the complete base
crossing and 45 the large-tier subset. Each cell is generated as
several independent instances, each instance is timed several times, and every
solver in the panel sees the same instance. The solver \texttt{"auto"} selects is timed
as a solver of its own, so its number carries the probing pass the rules cost,
and it is read against the fastest named solver in the same cell.

\begin{table}

\caption{\label{tab:regimes}Each dispatch rule over the cells of the regime grid where it fired. Solver is what the rule selects. Fastest counts the cells where that solver was also the quickest of the panel. Median and worst are the time the dispatched solver took over the time the cell's fastest named solver took, so 1.00x is a rule that picked the best available solver. The large tier times a reduced panel that omits the special-purpose solvers, so a cell whose dispatched solver is not in it reads below 1.00x.}
\centering
\begin{tabular}[t]{llrrrr}
\toprule
Rule & Solver & Cells & Fastest & Median & Worst\\
\midrule
no earlier rule applies & jv & 144 & 59 & 1.04x & 9.95x\\
no cost scale to exploit & hk01 & 45 & 35 & 1.00x & 1.93x\\
\bottomrule
\end{tabular}
\end{table}

Over the 189 cells the dispatcher picked the fastest solver of the
panel in 94 of them, at a median 1.02x of the
cell's best time. A cell reads below \texttt{1.00x} for one of two reasons. The large
tier times a reduced panel, so on its 9 binary cells the
solver the dispatcher named was not among those it is read against; on the
36 cells of that rule whose panel does hold its solver the median
is 1.02x. Elsewhere the two sides are separate timings of
one solver and scatter by a few percent either way. The two rules that were
removed can be read off the same table. Over the 96 cells with
at least three columns per row, where the rectangular rule used to fire, the
default takes a median
1.02x; over the 96 cells with more than
half the entries forbidden, 1.14x. Neither is a regime the
default needs help in. Where it is furthest behind, at
9.95x, the winner is auction\_scaled on a
500 x 500 problem with heavily tied costs, which is a structure no
rule in the table reads and none was added for: fitting one to this grid would
score it on the measurements that chose it.

The full grid, every solver in every cell, is in the supplementary material.

\FloatBarrier

\subsection{Comparison with established alternatives}\label{comparison-with-established-alternatives}

\subsubsection{Balance}\label{balance}

Speed only matters once the packages agree on the answer. We fix the task in
Figure \ref{fig:love}, one-to-one optimal
matching on a Mahalanobis distance with pooled within-group covariance, and run
it on the LaLonde NSW data \citep{LaLonde1986, DehejiaWahba1999}, 185 treated units,
429 controls and eight covariates. The pooled within-group convention is the one
\texttt{optmatch::match\_on()} uses by default and the one couplr adopts.

\begin{figure}[H]

{\centering \includegraphics[width=1\linewidth,alt={A dot plot with eight covariates on the vertical axis and absolute standardised mean difference on the horizontal axis. For each covariate an open circle marks the value before matching and a filled square the value after matching, joined by a grey line. Every covariate moves left toward zero. Five of the eight land below the dashed 0.1 threshold, married and 1974 earnings sit just above it near 0.12 and 0.13, and the indicator for Black respondents remains far to the right at 1.05.}]{rjournal_files/figure-latex/love-1} 

}

\caption{Absolute standardised mean differences on the eight LaLonde NSW covariates, before and after one-to-one optimal Mahalanobis matching with pooled within-group covariance. couplr, MatchIt and optmatch agree to three decimal places after matching, so a single matched point is shown per covariate. The dashed line marks the conventional 0.1 threshold. The indicator for Black respondents starts far outside the plotted range at 1.757 and remains at 1.053 after matching, a residual imbalance that no one-to-one matcher can resolve on this data.}\label{fig:love}
\end{figure}

The three packages produce identical absolute standardised mean differences to
three decimal places, so one matched point is plotted. Marginals could still
hide different assignments, so we score all three on one objective: the same
Mahalanobis distance matrix, summed over each package's chosen pairs. Each matches all 185 treated
units, at a total distance of 304.042 for couplr against 304.043 for both
alternatives, a relative difference of \(3 \times 10^{-6}\). That is agreement
on the objective to the precision reported here rather than a demonstration
that all three attain the same minimum, and separating those two readings is
what a certificate is for: couplr's matching comes back certified optimal for
this cost matrix, and the same check runs on either alternative's pairing
against the same duals.

Five of the eight covariates fall below 0.1 after matching, \texttt{married} and 1974
earnings sit just above at 0.124 and 0.128, and the indicator for Black
respondents drops from 1.757 to 1.053 without approaching it. That is a
property of this dataset rather than of any package: the control pool does not
hold enough Black respondents for a one-to-one match to balance that
indicator.

\FloatBarrier

\subsubsection{Scaling}\label{scaling}

Table \ref{tab:scaling} reports wall-clock time on synthetic problems with the
same eight-covariate structure, at six sizes with a one-to-two ratio of treated
to control units, so every problem here is rectangular. Every package in this
comparison materialises the distance matrix, so we time couplr with
\texttt{memory\_mode\ =\ "dense"}. Each size is generated as several independent
instances with the timing repetitions inside them, and the table reports the
median and the interquartile range across instances: repetitions of one matrix
measure the clock, and what varies between problems is what a reader wants to
know.

\begin{table}

\caption{\label{tab:scaling}One-to-one optimal Mahalanobis matching, wall-clock time by problem size. Treated to control ratio 1:2, eight covariates, pooled within-group covariance, single core with single-threaded BLAS on an Apple M4 Pro. Each cell is the median over independently generated instances of that size, with the interquartile range across instances in brackets and the timing repetitions taken inside each instance. The 50,000 row carries one instance and is a single run rather than a median, so no range is reported for it. The optmatch option max.problem.size was set to Inf for the two largest sizes. int overflow marks an integer-size overflow inside the optmatch back end reached through MatchIt; timeout marks a run exceeding the 600-second cap.}
\centering
\begin{tabular}[t]{lrrr}
\toprule
Problem size & couplr & optmatch & MatchIt\\
\midrule
167 + 333 & 17 ms (17-19) & 120 ms (111-121) & 134 ms (132-139)\\
667 + 1,333 & 140 ms (139-143) & 1.75 s (1.67-1.77) & 2.14 s (2.12-2.16)\\
1,667 + 3,333 & 789 ms (779-821) & 12.8 s (12.7-12.8) & 16.3 s (16.2-16.3)\\
3,333 + 6,667 & 3.16 s (3.1-3.22) & 60.5 s (59.8-60.5) & 74.2 s (73.5-74.3)\\
6,667 + 13,333 & 11.6 s (11.2-12.1) & 287 s (284-290) & int overflow\\
\addlinespace
16,667 + 33,333 & 76.3 s & timeout & timeout\\
\bottomrule
\end{tabular}
\end{table}

The margin widens with problem size. couplr is about
7 times faster than \CRANpkg{optmatch} at
\(n = 500\) and 25 times faster at
\(n = 20{,}000\), where it finishes in
11.6 seconds against
4.8 minutes.
\CRANpkg{MatchIt} runs
8
to
23
times slower up to \(n = 10{,}000\), which is the largest size it completes before
\texttt{matchit(method\ =\ "optimal")} aborts inside the \CRANpkg{optmatch} back end with
an integer-size overflow. At \(n = 50{,}000\) couplr finishes in
76
seconds and both alternatives exceed the 600-second cap.

What the table compares is three workflows on one task, from data frame to
matched pairs, not three implementations of one algorithm. The packages do not
state the same problem: \CRANpkg{optmatch} solves a minimum-cost flow problem,
which the full and variable-ratio designs it is built around need and a plain
one-to-one problem does not, and \CRANpkg{MatchIt} reaches that same back end,
while couplr's dispatcher sends this problem to Jonker-Volgenant. Reading the
margin as a faster inner loop would need a benchmark holding the formulation
fixed across packages, which this one does not.

The lazy path removes the memory ceiling as well. The same \(n = 50{,}000\) match
under \texttt{memory\_mode\ =\ "lazy"} completes in
62 seconds without ever
building the distance matrix, returning the assignment the dense path returns.
At \(n = 20{,}000\) it takes
6.4 seconds against
11.6
for the dense path, so on these problems recomputing distances is cheaper than
materialising and traversing the matrix.

\FloatBarrier

\subsubsection{Feature coverage}\label{feature-coverage}

Table \ref{tab:capability} lists where the three packages differ.
Rows where all three offer the feature, covariate distance, propensity-score
matching, exact blocking and a cost-matrix entry point, are omitted.

\begin{table}

\caption{\label{tab:capability}Selected assignment-layer features, omitting areas where the alternatives lead, among them the breadth of designs reachable through a single MatchIt call and the maturity of optmatch's full matching. optmatch's calipers threshold one distance object at a time, so a per-variable set requires combining objects, All three handle unequal group sizes; the second row records the formulation each reaches the solver with, couplr a rectangular cost matrix directly, optmatch's pairmatch() through fullmatch(), and MatchIt through optmatch::fullmatch(). The verifier row records whether a package exports dual potentials together with a function checking the optimality conditions against a returned solution. optmatch reports a gap instead: the exceedances attribute of a fullmatch() result is documented as an upper bound, not necessarily sharp, on how far the returned sum of distances exceeds the least possible sum over feasible solutions. Its node prices travel in an MCFSolutions attribute fullmatch()'s documentation does not describe, and the function scoring a primal against them is not exported. That row was assessed on 2026-08-26 and the others on 2026-08-09, both against MatchIt 4.7.2 and optmatch 0.10.8.}
\centering
\begin{tabular}[t]{lccc}
\toprule
Feature & couplr & MatchIt & optmatch\\
\midrule
Per-variable calipers from a named vector & yes & yes & partial\\
Rectangular cost matrix as solver entry point & yes & full match & full match\\
Sparse cost support & yes & no & yes\\
k-best assignments & yes & no & no\\
Bottleneck (minimax) assignment & yes & no & no\\
\addlinespace
Rosenbaum sensitivity bounds & yes & no & no\\
Public dual-potential verifier & yes & no & gap bound\\
User-selectable assignment algorithm & 19 solvers & via optmatch & 5 algorithms\\
\bottomrule
\end{tabular}
\end{table}

\FloatBarrier

\section{Summary and discussion}\label{summary-and-discussion}

We built couplr on the observation that one optimisation problem sits
underneath a set of tasks R currently treats as unrelated. Making the
assignment problem the organising object, rather than a hidden step inside a
matching function, is what lets one package serve a causal-inference user
starting from covariates and an operations-research user starting from a cost
matrix. Three consequences follow. Every flow-representable design compiles
into one flow model, so node potentials come back from all of them; those
potentials turn a reported status into a certificate; and a certificate is what
makes it safe to solve a problem that was never built, since a loop growing a
sparse graph needs a statement about the complete one to know when to stop.

What that architecture buys is measurable. At 16,667 + 33,333 the restricted graph
settled at 0.29\% of the complete pair count, returned a certified
optimum, and ran 3.1
times faster than traversing the complete graph lazily. A sweep of
20 caliper values solved as one warm-started path cost a median
0.47 of solving each point cold, and
across the 189 cells of the solver-regime grid the dispatcher ran
at a median 1.02x of the cell's fastest measured solver.

couplr has five limitations:

\begin{itemize}
\tightlist
\item
  The edge-generation loop's scope is the unit-capacity assignment: a full
  match's column nodes carry capacities above one, so its column duals are not
  the duals the pricer reads and \texttt{full\_match()} has no implicit path. That is a
  limit of this implementation rather than of the method, since a general flow
  arc has a reduced cost from the potentials at its ends.
\item
  The lazy cost path covers one-to-one matching without replacement under
  Jonker-Volgenant and auction with the built-in metrics. Replacement, a ratio
  above one, \texttt{cardinality\_match()} and a user-supplied distance function, which
  expects one call over both matrices rather than one per pair, build the dense
  matrix.
\item
  \texttt{cardinality\_match()} answers a moment-constrained problem by branch and
  bound under a node and time limit, so a run can stop with a gap rather than a
  certificate and reports which it reached. Fine and refined balance, which the
  network represents directly, is answered certified.
\item
  The dispatcher reads two properties and otherwise defaults to
  Jonker-Volgenant, which the regime grid supports across most of the
  structures it covers. Over the small square cells with most of the matrix
  forbidden, scaled auction is quicker in most of them, at a median of about
  three times and as much as ten where the costs are also heavily tied. Fitting
  a rule to this grid would score it on the measurements that chose it, so the
  default stands and every solver stays reachable by name.
\item
  Matching on observed covariates cannot address unmeasured confounding, which
  the sensitivity analysis quantifies rather than removes.
\end{itemize}

Two extensions follow. Pricing a general flow arc from the potentials at the
nodes it joins would carry the loop from the unit-capacity assignment to every
design the flow model compiles, and with it to \texttt{full\_match()}. A bound tighter
than a ball would extend the pricer past the dimension it stops paying at.

couplr is available on CRAN and developed at
\texttt{https://github.com/gcol33/couplr}, with documentation at
\texttt{https://gillescolling.com/couplr/}. Bug reports and feature requests go
through the GitHub issue tracker. No external funding supported this work.

\bibliography{RJreferences.bib}

\address{%
Gilles Colling\\
University of Vienna\\%
Department of Botany and Biodiversity Research\\ Rennweg 14, 1030 Vienna, Austria\\
%
\url{https://gillescolling.com}\\%
\textit{ORCID: \href{https://orcid.org/0000-0003-3070-6066}{0000-0003-3070-6066}}\\%
\href{mailto:gilles.colling051@gmail.com}{\nolinkurl{gilles.colling051@gmail.com}}%
}
